\documentclass[a4paper,10pt,brazil]{article}
\usepackage{provastyle}
\usepackage{menukeys}

\begin{document}
\makeexamtitlepage{UFOP}{BCC222: Programação Funcional}{José Romildo}{2026/1}{Primeira Prova}{07}
\begin{multicols}{2}
\questionTitle{1}{Algoritmo de Euclides}
\begin{flushright}
  \footnotesize
  \textit{\textbf{9: Funções recursivas}}\\
  \textit{c09-projeto-mdc-01}\\

\end{flushright}
Qual alternativa expressa corretamente o passo recursivo do algoritmo de Euclides para calcular \haskellinline|mdc a b| quando \haskellinline|b > 0|?

\begin{enumerate}
  \item Retornar imediatamente \haskellinline|a + b|.

  \item Chamar recursivamente \haskellinline|mdc (a + b) b|.

  \item Chamar recursivamente \haskellinline|mdc (a - b) (b - a)| sem verificar sinais.

  \item Chamar recursivamente \haskellinline|mdc b (mod a b)|.

  \item Chamar recursivamente \haskellinline|mdc a (mod b a)|.


\end{enumerate}

\questionTitle{2}{Refatoração para forma mais concisa}
\begin{flushright}
  \footnotesize
  \textit{\textbf{5: Expressão condicional}}\\
  \textit{c05-if-vs-guardas-02}\\

\end{flushright}
Considere a função:
\begin{haskellcode}
classifica n
  | n > 0     = "Positivo"
  | n < 0     = "Negativo"
  | n == 0    = "Zero"
  | otherwise = "Impossível"
\end{haskellcode}
Qual definição é funcionalmente equivalente, mas mais concisa?

\begin{enumerate}
  \item \begin{haskellcode}
classifica n
  | n > 0     = "Positivo"
  | otherwise = "Negativo"
\end{haskellcode}

  \item \begin{haskellcode}
classifica n
  | n >= 0    = "Positivo"
  | otherwise = "Negativo"
\end{haskellcode}

  \item \begin{haskellcode}
classifica n
  | n == 0    = "Zero"
  | otherwise = "Positivo ou Negativo"
\end{haskellcode}

  \item A função original não pode ser simplificada sem mudar seu comportamento.

  \item \begin{haskellcode}
classifica n
  | n > 0     = "Positivo"
  | n < 0     = "Negativo"
  | otherwise = "Zero"
\end{haskellcode}


\end{enumerate}

\questionTitle{3}{Modelagem de inventário}
\begin{flushright}
  \footnotesize
  \textit{\textbf{6: Estruturas de dados básicas}}\\
  \textit{c06-modelagem-02}\\

\end{flushright}
Você quer representar um inventário como uma lista de produtos. Cada produto possui nome (\haskellinline{String}), código (\haskellinline{Int}) e preço opcional (\haskellinline{Maybe Double}). Qual tipo representa corretamente esse inventário?

\begin{enumerate}
  \item \haskellinline{[String, Int, Maybe Double]}

  \item \haskellinline{Maybe (String, Int, Double)}

  \item \haskellinline{[(String, Int, Maybe Double)]}

  \item \haskellinline{(String, Int, Maybe Double)}

  \item \haskellinline{([String], [Int], [Maybe Double])}


\end{enumerate}

\questionTitle{4}{Prompt com flush em programa completo}
\begin{flushright}
  \footnotesize
  \textit{\textbf{8: Programas interativos}}\\
  \textit{c08-programas-completos-05}\\

\end{flushright}
Considere o trecho:

\begin{minted}{haskell}
main :: IO ()
main = do
  putStr "Palavra: "
  hFlush stdout
  palavra <- getLine
  putStrLn palavra
\end{minted}

Qual é a finalidade de \haskellinline|hFlush stdout| nesse programa?

\begin{enumerate}
  \item Adicionar uma quebra de linha ao final do \emph{prompt}.

  \item Imprimir automaticamente o conteúdo de \haskellinline|palavra|.

  \item Garantir que o \emph{prompt} produzido por \haskellinline|putStr| seja exibido antes de o programa aguardar a entrada com \haskellinline|getLine|.

  \item Limpar a entrada digitada pelo usuário antes de \haskellinline|getLine|.

  \item Converter \haskellinline|palavra| de \haskellinline|IO String| para \haskellinline|String|.


\end{enumerate}

\questionTitle{5}{Interpretação de Num a => a}
\begin{flushright}
  \footnotesize
  \textit{\textbf{4: Tipos de Dados}}\\
  \textit{c04-ghci-interpretacao-01}\\

\end{flushright}
Ao executar \texttt{:type 42} no GHCi, a resposta é \haskellinline{42 :: Num a => a}. O que isso significa?

\begin{enumerate}
  \item Indica um erro de tipo, pois o tipo é ambíguo e não pode ser resolvido.

  \item \haskellinline{42} é uma constante numérica que só pode ser do tipo \haskellinline{Int}.

  \item \haskellinline{42} pode ser interpretado como um valor de qualquer tipo que pertença à classe \haskellinline{Num} (como \haskellinline{Int}, \haskellinline{Integer}, \haskellinline{Double}).

  \item \haskellinline{42} é um valor de um tipo desconhecido que não pode ser usado em operações aritméticas.

  \item Significa que \haskellinline{42} é uma função que recebe um \haskellinline{Num} e retorna um \haskellinline{a}.


\end{enumerate}

\questionTitle{6}{Where não é expressão}
\begin{flushright}
  \footnotesize
  \textit{\textbf{3: Expressões e definições}}\\
  \textit{c03-definicoes-where-04}\\

\end{flushright}
Por que a forma abaixo não é uma maneira válida de criar uma definição local em Haskell?

\begin{haskellcode}
resultado = where x = 1 in x + 1
\end{haskellcode}

\begin{enumerate}
  \item Porque nomes definidos por \haskellinline|where| precisam começar com letra maiúscula.

  \item Porque \haskellinline|where| não é uma expressão do tipo \haskellinline|where ... in ...|; ele é uma cláusula anexada a uma equação.

  \item Porque \haskellinline|where| só pode ser usado dentro do GHCi, nunca em arquivos fonte.

  \item Porque \haskellinline|where| não permite definir valores numéricos.

  \item Porque \haskellinline|x + 1| deveria ser escrito obrigatoriamente como \haskellinline|(+ x 1)|.


\end{enumerate}

\questionTitle{7}{Inferência de tipo de isDigit}
\begin{flushright}
  \footnotesize
  \textit{\textbf{4: Tipos de Dados}}\\
  \textit{c04-inferencia-basica-02}\\

\end{flushright}
Considere \haskellinline{import Data.Char}. Qual é o tipo inferido para a expressão \haskellinline{isDigit '5'}?

\begin{enumerate}
  \item \haskellinline{Char -> Bool}

  \item \haskellinline{True}

  \item \haskellinline{Char}

  \item \haskellinline{Bool}

  \item \haskellinline{Num a => a}


\end{enumerate}

\questionTitle{8}{Return não encerra um bloco do}
\begin{flushright}
  \footnotesize
  \textit{\textbf{10: Ações de E/S recursivas}}\\
  \textit{c10-return-tipos-05}\\

\end{flushright}
Muitos alunos com experiência em linguagens imperativas associam \haskellinline|return| ao encerramento imediato de uma função. Considere:

\begin{haskellcode}
interacao :: IO ()
interacao = do
  putStrLn "Início"
  return "valor intermediário"
  putStrLn "Fim"
\end{haskellcode}

Qual alternativa descreve corretamente o comportamento desse bloco?

\begin{enumerate}
  \item O programa não compila, pois \haskellinline|return| só pode aparecer na última linha de um bloco \haskellinline|do|.

  \item A execução para em \haskellinline|return "valor intermediário"|, portanto \haskellinline|"Fim"| não é impresso.

  \item O programa imprime \haskellinline|"valor intermediário"| automaticamente entre as duas mensagens.

  \item A ação \haskellinline|return| extrai uma string de dentro de \haskellinline|IO|, sendo o inverso de \haskellinline|<-|.

  \item A mensagem \haskellinline|"Início"| é impressa, a ação \haskellinline|return "valor intermediário"| produz uma string que é descartada, e a mensagem \haskellinline|"Fim"| também é impressa.


\end{enumerate}

\questionTitle{9}{Abreviações sintáticas de listas e strings}
\begin{flushright}
  \footnotesize
  \textit{\textbf{3: Expressões e definições}}\\
  \textit{c03-comentarios-valores-04}\\

\end{flushright}
Qual alternativa apresenta uma equivalência correta de açúcar sintático em Haskell?

\begin{enumerate}
  \item \haskellinline|"abc"| é uma abreviação para \haskellinline|['a','b','c']|.

  \item \haskellinline|[]| é uma abreviação para a tupla vazia \haskellinline|()|.

  \item \haskellinline|[1,2]| é uma abreviação para a tupla \haskellinline|(1,2)|.

  \item \haskellinline|"abc"| é uma abreviação para \haskellinline|["a","b","c"]|.

  \item \haskellinline|(1,2,3)| é uma abreviação para a lista \haskellinline|[1,2,3]|.


\end{enumerate}

\questionTitle{10}{Recursão e funções de redução}
\begin{flushright}
  \footnotesize
  \textit{\textbf{9: Funções recursivas}}\\
  \textit{c09-eficiencia-reducao-01}\\

\end{flushright}
O capítulo observa que, em linguagens funcionais, repetições podem ser expressas por recursão ou por abstrações construídas sobre padrões recursivos. Qual alternativa expressa melhor essa ideia?

\begin{enumerate}
  \item Funções de redução, como \haskellinline|foldr| e \haskellinline|foldl|, encapsulam padrões recorrentes de recursão sobre estruturas.

  \item Funções de redução só existem em linguagens imperativas.

  \item Funções de redução impedem o uso de funções puras.

  \item Funções de redução são equivalentes a variáveis globais atualizadas em laços.

  \item Funções de redução substituem completamente a necessidade de compreender casos base.


\end{enumerate}

\questionTitle{11}{Primeira guarda verdadeira vence}
\begin{flushright}
  \footnotesize
  \textit{\textbf{5: Expressão condicional}}\\
  \textit{c05-guardas-avaliacao-02}\\

\end{flushright}
Considere a função:
\begin{haskellcode}
f x
  | x >= 0    = 10
  | x > 3     = 20
  | otherwise = 30
\end{haskellcode}
Qual é o valor de \haskellinline|f 4|?

\begin{enumerate}
  \item \haskellinline|30|

  \item Ocorre um erro, pois duas guardas são verdadeiras.

  \item \haskellinline|20|

  \item \haskellinline|10|

  \item O compilador rejeita a definição pela sobreposição das guardas.


\end{enumerate}

\questionTitle{12}{Avaliação preguiçosa e acumuladores}
\begin{flushright}
  \footnotesize
  \textit{\textbf{10: Ações de E/S recursivas}}\\
  \textit{c10-acumuladores-cauda-05}\\

\end{flushright}
No capítulo, a versão com acumulador é apresentada como uma melhoria estrutural em relação à versão que deixa somas pendentes. Em Haskell, qual ressalva técnica é importante ao discutir uso de memória nessa versão?

\begin{enumerate}
  \item A ressalva é que \haskellinline|readLn| avalia todos os valores futuros antes de retornar o primeiro.

  \item A ressalva é que \haskellinline|return| sempre força completamente o valor acumulado antes de empacotá-lo em \haskellinline|IO|.

  \item A ressalva é que toda função recursiva de cauda em Haskell é rejeitada pelo compilador por não ser declarativa.

  \item A ressalva é que acumuladores não podem ser usados em ações de E/S, apenas em funções puras.

  \item Mesmo com chamada em posição de cauda, a avaliação preguiçosa pode acumular expressões não avaliadas no acumulador se ele não for forçado adequadamente.


\end{enumerate}

\questionTitle{13}{Responsabilidades no fechamento de notas}
\begin{flushright}
  \footnotesize
  \textit{\textbf{10: Ações de E/S recursivas}}\\
  \textit{c10-separacao-puro-04}\\

\end{flushright}
Em um programa de fechamento de notas, uma boa organização separa ações de E/S e processamento puro. Qual distribuição de responsabilidades é mais adequada?

\begin{enumerate}
  \item Colocar leitura, cálculo de nota final, classificação e impressão em uma única função pura, sem tipo \haskellinline|IO|.

  \item Evitar listas, pois programas com \haskellinline|IO| não podem armazenar dados lidos para processamento posterior.

  \item Usar ações de E/S para ler dados e exibir resultados, deixando cálculos como média, nota final, situação e percentuais em funções puras.

  \item Fazer com que cada função pura leia novamente a matrícula e as notas do aluno para evitar passagem de parâmetros.

  \item Usar \haskellinline|return| para converter todas as funções puras em ações \haskellinline|IO|, mesmo quando não há interação externa.


\end{enumerate}

\questionTitle{14}{Recursão para percorrer uma lista com efeitos}
\begin{flushright}
  \footnotesize
  \textit{\textbf{10: Ações de E/S recursivas}}\\
  \textit{c10-lacos-io-05}\\

\end{flushright}
Considere novamente:

\begin{haskellcode}
mostraLista :: Show a => [a] -> IO ()
mostraLista lista
  | null lista = return ()
  | otherwise = do
      print (head lista)
      mostraLista (tail lista)
\end{haskellcode}

Qual alternativa explica corretamente a repetição realizada por essa ação?

\begin{enumerate}
  \item A lista é percorrida por \haskellinline|map| implicitamente, pois \haskellinline|IO| é uma mônada.

  \item A repetição ocorre porque \haskellinline|print| chama automaticamente \haskellinline|mostraLista| para cada elemento da lista.

  \item A repetição ocorre porque \haskellinline|null lista| percorre a lista inteira e imprime cada elemento.

  \item Em cada passo, a ação imprime o primeiro elemento e chama recursivamente \haskellinline|mostraLista| com a cauda da lista.

  \item A função só imprime o primeiro elemento, pois a chamada recursiva é descartada por ter tipo \haskellinline|IO ()|.


\end{enumerate}

\questionTitle{15}{Escopo limitado à equação}
\begin{flushright}
  \footnotesize
  \textit{\textbf{5: Expressão condicional}}\\
  \textit{c05-where-escopo-02}\\

\end{flushright}
Considere a definição:
\begin{haskellcode}
f x
  | y > 0     = y
  | otherwise = 0
  where y = x - 1

f x = y + 10
\end{haskellcode}
Qual afirmação está correta?

\begin{enumerate}
  \item Ambas as equações compartilham \haskellinline|y|, mas apenas se houver uma única assinatura de tipo.

  \item A segunda equação não pode usar \haskellinline|y|, pois esse nome vale apenas na primeira equação.

  \item A segunda equação pode usar \haskellinline|y| normalmente, pois \haskellinline|where| tem escopo em toda a função.

  \item A segunda equação pode usar \haskellinline|y|, mas somente se ela também tiver guardas.

  \item \haskellinline|where| sempre define variáveis globais no módulo.


\end{enumerate}

\questionTitle{16}{Soma com sucessor e antecessor}
\begin{flushright}
  \footnotesize
  \textit{\textbf{9: Funções recursivas}}\\
  \textit{c09-rastreamento-soma-01}\\

\end{flushright}
Considere:
\begin{haskellcode}
soma :: Integer -> Integer -> Integer
soma m n
  | n == 0    = m
  | otherwise = soma (succ m) (pred n)
\end{haskellcode}
Qual é o valor de \haskellinline|soma 9 1|?

\begin{enumerate}
  \item \haskellinline|1|

  \item \haskellinline|9|

  \item \haskellinline|11|

  \item \haskellinline|8|

  \item \haskellinline|10|


\end{enumerate}

\questionTitle{17}{Ausência de segurança estrita em sinônimos}
\begin{flushright}
  \footnotesize
  \textit{\textbf{4: Tipos de Dados}}\\
  \textit{c04-sinonimos-equivalencia-01}\\

\end{flushright}
Considere:
\begin{haskellcode}
type Nome = String
type Endereco = String
geraEtiqueta :: Nome -> Endereco -> String
geraEtiqueta n e = "Para: " ++ n ++ " - Env: " ++ e
\end{haskellcode}
Se o programador chamar \haskellinline{geraEtiqueta "Rua X" "João"} (invertendo os argumentos), o que acontece?

\begin{enumerate}
  \item O código compila, mas lança uma exceção em tempo de execução.

  \item O código compila sem erros ou avisos, pois \haskellinline{Nome} e \haskellinline{Endereco} são o mesmo tipo que \haskellinline{String}.

  \item O compilador rejeita o código com um erro de tipo, protegendo contra a inversão.

  \item Ocorre um erro sintático porque sinônimos exigem construtores (ex: \haskellinline{Nome "João"}).

  \item O compilador emite um aviso, mas gera o executável.


\end{enumerate}

\questionTitle{18}{Ordem de avaliação das guardas}
\begin{flushright}
  \footnotesize
  \textit{\textbf{5: Expressão condicional}}\\
  \textit{c05-guardas-sintaxe-semantica-02}\\

\end{flushright}
Como as guardas de uma definição são avaliadas?

\begin{enumerate}
  \item O compilador reordena as guardas para testar primeiro as mais específicas.

  \item Todas as guardas são avaliadas, e a última verdadeira determina o resultado.

  \item De cima para baixo; a primeira guarda verdadeira determina o resultado.

  \item As guardas são avaliadas de baixo para cima.

  \item O resultado é escolhido pela guarda cuja expressão à direita de \haskellinline|=| parece mais simples.


\end{enumerate}

\questionTitle{19}{Escolha entre arredondar e converter integralmente}
\begin{flushright}
  \footnotesize
  \textit{\textbf{7: Polimorfismo}}\\
  \textit{cap07-conversao-arredondamento-escolha}\\

\end{flushright}
Uma função deve receber um \haskellinline{Double} e produzir um \haskellinline{Int} arredondando para o inteiro mais próximo. Qual função é adequada?
\begin{enumerate}
  \item \haskellinline{realToFrac}, pois sempre retorna um tipo integral
  \item \haskellinline{fromIntegral}
  \item \haskellinline{show}, pois a conversão passa por texto
  \item \haskellinline{round}
  \item \haskellinline{fromInteger}

\end{enumerate}

\questionTitle{20}{Acesso composto a pares}
\begin{flushright}
  \footnotesize
  \textit{\textbf{6: Estruturas de dados básicas}}\\
  \textit{c06-tuplas-acesso-01}\\

\end{flushright}
Considere a expressão abaixo.
\begin{minted}{haskell}
fst (snd (False, ("Bruno", 34)))
\end{minted}
Qual é o seu resultado?

\begin{enumerate}
  \item \haskellinline{ 34 }

  \item \haskellinline{ ("Bruno",34) }

  \item \haskellinline{ False }

  \item \haskellinline{ "Bruno" }

  \item A expressão resulta em erro de compilação.


\end{enumerate}

\questionTitle{21}{Lookup com valor padrão}
\begin{flushright}
  \footnotesize
  \textit{\textbf{6: Estruturas de dados básicas}}\\
  \textit{c06-lookup-associacao-02}\\

\end{flushright}
Qual é o resultado da expressão abaixo?
\begin{minted}{haskell}
let tabela = [("alice",1.0),("bob",2.0)]
in fromMaybe 0.0 (lookup "carlos" tabela)
\end{minted}

\begin{enumerate}
  \item \haskellinline{ 2.0 }

  \item \haskellinline{ 0.0 }

  \item \haskellinline{ Just 0.0 }

  \item Ocorre um erro em tempo de execução.

  \item \haskellinline{ Nothing }


\end{enumerate}

\questionTitle{22}{Aninhamento com três casos}
\begin{flushright}
  \footnotesize
  \textit{\textbf{5: Expressão condicional}}\\
  \textit{c05-if-aninhamento-02}\\

\end{flushright}
Avalie a função abaixo na entrada \haskellinline|f 5|:
\begin{haskellcode}
f x = if x > 0
      then if x > 10
           then 1
           else 2
      else 3
\end{haskellcode}

\begin{enumerate}
  \item \haskellinline|3|

  \item \haskellinline|2|

  \item \haskellinline|1|

  \item Ocorre um erro de sintaxe por falta de parênteses.

  \item Ocorre um erro de ambiguidade entre os dois \haskellinline|else|.


\end{enumerate}

\questionTitle{23}{Análise de restrições combinadas}
\begin{flushright}
  \footnotesize
  \textit{\textbf{7: Polimorfismo}}\\
  \textit{cap07-expressao-comparacao-show}\\

\end{flushright}
Considere:

\begin{minted}{haskell}
h x y = x <= y && show x == show y
\end{minted}

Qual é a assinatura mais geral de \haskellinline{h}?

\begin{enumerate}
  \item \mintinline{haskell}{h :: (Ord a, Show a) => a -> a -> Bool}
  \item \mintinline{haskell}{h :: Ord a => a -> a -> String}
  \item \mintinline{haskell}{h :: a -> a -> Bool}
  \item \mintinline{haskell}{h :: Show a => a -> a -> String}
  \item \mintinline{haskell}{h :: Eq a => a -> a -> Bool}

\end{enumerate}

\questionTitle{24}{Regra de layout em do}
\begin{flushright}
  \footnotesize
  \textit{\textbf{8: Programas interativos}}\\
  \textit{c08-blocos-do-05}\\

\end{flushright}
Em um bloco \haskellinline|do| escrito sem chaves e ponto e vírgula explícitos, qual regra de layout é essencial?

\begin{enumerate}
  \item Cada linha deve começar obrigatoriamente na primeira coluna do arquivo.

  \item A indentação é ignorada em Haskell; apenas quebras de linha importam.

  \item A primeira ação do bloco deve ficar menos indentada que a palavra \haskellinline|do|.

  \item As ações devem ficar cada vez mais indentadas à medida que aparecem no bloco.

  \item As ações pertencentes ao mesmo bloco devem ficar alinhadas na mesma coluna de indentação, a partir do primeiro item do bloco.


\end{enumerate}

\questionTitle{25}{Entrada que provoca falha}
\begin{flushright}
  \footnotesize
  \textit{\textbf{5: Expressão condicional}}\\
  \textit{c05-guardas-exaustividade-02}\\

\end{flushright}
Para qual chamada a função abaixo falha em tempo de execução?
\begin{haskellcode}
comparar x y
  | x > y = "Maior"
  | x < y = "Menor"
\end{haskellcode}

\begin{enumerate}
  \item \haskellinline|comparar 5 5|

  \item \haskellinline|comparar 0 (-1)|

  \item \haskellinline|comparar 10 5|

  \item Nenhuma, pois o compilador adiciona \haskellinline|otherwise| implicitamente.

  \item \haskellinline|comparar 5 10|


\end{enumerate}

\questionTitle{26}{A primeira linguagem funcional prática}
\begin{flushright}
  \footnotesize
  \textit{\textbf{1: Paradigmas de programação}}\\
  \textit{c01-historia-evolucao-02}\\

\end{flushright}
Criada na virada da década de 1950 por John McCarthy no MIT, qual linguagem de programação é historicamente consagrada como o primeiro esforço prático para incorporar o paradigma funcional, introduzindo conceitos inovadores à época, como a recursão nativa e o \emph{Garbage Collection}?

\begin{enumerate}
  \item ML (MetaLanguage)

  \item Miranda

  \item Fortran

  \item Lisp

  \item COBOL


\end{enumerate}

\questionTitle{27}{If como subexpressão}
\begin{flushright}
  \footnotesize
  \textit{\textbf{5: Expressão condicional}}\\
  \textit{c05-if-avaliacao-02}\\

\end{flushright}
Considere a expressão:
\begin{haskellcode}
4 * (if 'B' < 'A' then 2 + 3 else 2) - 3
\end{haskellcode}
Qual é o seu valor?

\begin{enumerate}
  \item Ocorre um erro de tipos.

  \item \haskellinline|20|

  \item \haskellinline|-1|

  \item \haskellinline|17|

  \item \haskellinline|5|


\end{enumerate}

\questionTitle{28}{Classe mínima para div e mod}
\begin{flushright}
  \footnotesize
  \textit{\textbf{7: Polimorfismo}}\\
  \textit{cap07-hierarquia-classe-minima-divisao}\\

\end{flushright}
Considere uma função que usa apenas \haskellinline{div} e \haskellinline{mod} sobre seus argumentos. Qual restrição é a mais adequada para a assinatura mais geral?
\begin{enumerate}
  \item \haskellinline{Fractional a}, pois divisão sempre exige tipos fracionários.
  \item \haskellinline{Integral a}, pois \haskellinline{div} e \haskellinline{mod} operam sobre tipos integrais.
  \item \haskellinline{Floating a}, pois qualquer divisão exige ponto flutuante.
  \item Nenhuma restrição é necessária, pois \haskellinline{div} é paramétrico puro.
  \item \haskellinline{Show a}, pois \haskellinline{div} converte os operandos para texto.

\end{enumerate}

\questionTitle{29}{Valor puro e ação de E/S}
\begin{flushright}
  \footnotesize
  \textit{\textbf{8: Programas interativos}}\\
  \textit{c08-tipos-acoes-05}\\

\end{flushright}
Qual alternativa descreve corretamente a diferença entre \haskellinline|"abc"| e \haskellinline|getLine|?

\begin{enumerate}
  \item Ambas têm tipo \haskellinline|IO String|, pois toda string em Haskell pode ser impressa no terminal.

  \item \haskellinline|"abc"| tem tipo \haskellinline|IO ()| e \haskellinline|getLine| tem tipo \haskellinline|String|.

  \item Ambas têm tipo \haskellinline|String|, mas \haskellinline|getLine| é avaliada de forma preguiçosa.

  \item \haskellinline|getLine| é uma função pura do tipo \haskellinline|() -> String|, equivalente a uma função sem argumentos.

  \item \haskellinline|"abc"| é uma \haskellinline|String| pura; \haskellinline|getLine| é uma ação do tipo \haskellinline|IO String| que pode produzir uma \haskellinline|String| quando executada.


\end{enumerate}

\questionTitle{30}{Tipo da condição}
\begin{flushright}
  \footnotesize
  \textit{\textbf{5: Expressão condicional}}\\
  \textit{c05-if-regras-02}\\

\end{flushright}
Qual das expressões abaixo falha especificamente porque a condição do \haskellinline|if| não tem tipo \haskellinline|Bool|?

\begin{enumerate}
  \item \haskellinline|if True then 10 else 20|

  \item \haskellinline|if null [] then 10 else 20|

  \item \haskellinline|if 3 < 5 then 10 else 20|

  \item \haskellinline|if even 4 then 10 else 20|

  \item \haskellinline|if 3 then 10 else 20|


\end{enumerate}

\questionTitle{31}{Composição com fromMaybe}
\begin{flushright}
  \footnotesize
  \textit{\textbf{6: Estruturas de dados básicas}}\\
  \textit{c06-maybe-operacoes-01}\\

\end{flushright}
Qual é o resultado da expressão abaixo?
\begin{minted}{haskell}
fromMaybe "sem total" (if False then Just "calculado" else Nothing)
\end{minted}

\begin{enumerate}
  \item \haskellinline{ Just "calculado" }

  \item \haskellinline{ "calculado" }

  \item Ocorre um erro de tipos incompatíveis.

  \item \haskellinline{ Nothing }

  \item \haskellinline{ "sem total" }


\end{enumerate}

\questionTitle{32}{Análise declarativa da abstração do Quicksort}
\begin{flushright}
  \footnotesize
  \textit{\textbf{1: Paradigmas de programação}}\\
  \textit{c01-mecanismo-reducao-02}\\

\end{flushright}
Considere a cláusula principal da implementação declarativa do \emph{Quicksort} em Haskell:

\begin{haskellcode}
qsort (pivo:resto) = qsort (filter (<= pivo) resto) ++ 
                     [pivo] ++ 
                     qsort (filter (> pivo) resto)
\end{haskellcode}

Ao comparar este trecho com a tradicional implementação iterativa em C, qual das seguintes afirmações avalia corretamente a diferença de abstração utilizada na partição dos elementos?

\begin{enumerate}
  \item A versão funcional oculta a passagem de ponteiros de memória nos bastidores do operador \haskellinline|:|, realizando uma mutação de array idêntica à troca de variáveis (\cinline|swap|) do algoritmo em C.

  \item A versão funcional abstrai a manipulação de índices de memória e laços \cinline|while|, descrevendo a lista particionada como o resultado da aplicação de funções de filtragem baseadas em propriedades lógicas da lista.

  \item A versão funcional falha em implementar a estratégia de \emph{divisão e conquista}, pois a linguagem Haskell não suporta chamadas de recursão de cauda nativa como a linguagem C.

  \item A versão funcional requer que a lista de entrada seja previamente ordenada, caso contrário a função \haskellinline|filter| lançará uma exceção de limite de tempo de execução.

  \item A versão funcional é inerentemente mais ineficiente pois exige que o compilador traduza o operador de concatenação (\haskellinline|++|) em loops duplamente encadeados estritos.


\end{enumerate}

\questionTitle{33}{Hierarquia de Bibliotecas: Hackage vs Prelude}
\begin{flushright}
  \footnotesize
  \textit{\textbf{2: Ambiente de desenvolvimento}}\\
  \textit{c02-ecossistema-ferramentas-04}\\

\end{flushright}
O ecossistema Haskell divide as bibliotecas e módulos em diferentes níveis de acessibilidade. Em relação a essa estrutura, qual é a principal diferença prática entre o \textbf{Prelude} e o repositório \textbf{Hackage}?

\begin{enumerate}
  \item O Prelude é utilizado exclusivamente para compilar programas gráficos de janela, enquanto o Hackage é restrito ao desenvolvimento de aplicações de linha de comando.

  \item O Prelude precisa ser importado manualmente em todo arquivo \texttt{.hs}, ao passo que todas as bibliotecas do Hackage já vêm instaladas e prontas para uso direto no GHCi.

  \item Não há diferença; ambos são nomes distintos para a mesma ferramenta interna do GHC responsável por otimizar a avaliação preguiçosa.

  \item O Prelude é carregado automaticamente pelo compilador contendo funções básicas, enquanto o Hackage é um repositório online de onde se baixa pacotes de terceiros usando ferramentas como o Cabal.

  \item O Prelude contém as bibliotecas fechadas pagas por licença comercial, enquanto o Hackage é o acervo restrito a componentes de código livre (\emph{open-source}).


\end{enumerate}

\questionTitle{34}{Assinatura mais geral e uso excessivo de restrições}
\begin{flushright}
  \footnotesize
  \textit{\textbf{7: Polimorfismo}}\\
  \textit{cap07-param-funcao-restrita-vs-geral}\\

\end{flushright}
Um estudante propôs a seguinte assinatura para a função abaixo:

\begin{minted}{haskell}
duplicaPar x = (x, x)

duplicaPar :: Eq a => a -> (a, a)
\end{minted}

Qual é a melhor avaliação técnica?

\begin{enumerate}
  \item A função deveria exigir \haskellinline{Ord a}, pois tuplas sempre precisam ser ordenáveis.
  \item A assinatura correta é \mintinline{haskell}{duplicaPar :: a -> a}, pois a tupla não altera o tipo do valor.
  \item A função é monomórfica, pois o mesmo \haskellinline{x} aparece duas vezes no corpo.
  \item A restrição \haskellinline{Eq a} é desnecessária; a assinatura mais geral é \mintinline{haskell}{duplicaPar :: a -> (a, a)}.
  \item A restrição \haskellinline{Eq a} é obrigatória, pois repetir um valor exige compará-lo consigo mesmo.

\end{enumerate}

\questionTitle{35}{Strings como listas}
\begin{flushright}
  \footnotesize
  \textit{\textbf{6: Estruturas de dados básicas}}\\
  \textit{c06-listas-operacoes-02}\\

\end{flushright}
Como \haskellinline{String} é uma lista de caracteres, a expressão abaixo pode ser avaliada pelas mesmas funções de listas.
\begin{minted}{haskell}
drop 3 "Funcional"
\end{minted}
Qual é o resultado?

\begin{enumerate}
  \item \haskellinline{ "Fun" }

  \item \haskellinline{ "cional" }

  \item \haskellinline{ "uncional" }

  \item \haskellinline{ "ional" }

  \item A expressão resulta em erro de tipo.


\end{enumerate}

\questionTitle{36}{Classe mínima para igualdade e ordenação}
\begin{flushright}
  \footnotesize
  \textit{\textbf{7: Polimorfismo}}\\
  \textit{cap07-adhoc-classe-minima-comparacao}\\

\end{flushright}
Considere a função:

\begin{minted}{haskell}
entre x y z = x <= y && y <= z
\end{minted}

Qual restrição de classe é suficiente e mais geral para o tipo dos argumentos?

\begin{enumerate}
  \item \haskellinline{Ord a}, pois \haskellinline{(<=)} pertence a \haskellinline{Ord} e \haskellinline{Ord} já pressupõe \haskellinline{Eq}.
  \item \haskellinline{Num a}, pois qualquer comparação exige que os valores sejam numéricos.
  \item \haskellinline{Eq a}, pois igualdade é suficiente para implementar \haskellinline{(<=)}.
  \item Nenhuma restrição é necessária, pois \haskellinline{(<=)} é paramétrico puro.
  \item \haskellinline{Show a}, pois o operador \haskellinline{(&&)} precisa imprimir os operandos.

\end{enumerate}

\questionTitle{37}{Regra de nomenclatura de tipos}
\begin{flushright}
  \footnotesize
  \textit{\textbf{4: Tipos de Dados}}\\
  \textit{c04-conceito-nomes-01}\\

\end{flushright}
Qual é a regra sintática obrigatória para nomes de tipos em Haskell?

\begin{enumerate}
  \item Não podem conter caracteres especiais.

  \item Devem ser escritos inteiramente em maiúsculas.

  \item Devem conter pelo menos um dígito.

  \item Devem começar com uma letra maiúscula.

  \item Devem começar com uma letra minúscula.


\end{enumerate}

\questionTitle{38}{Duas formas de construir uma lista lida}
\begin{flushright}
  \footnotesize
  \textit{\textbf{10: Ações de E/S recursivas}}\\
  \textit{c10-listas-ordem-04}\\

\end{flushright}
Compare as duas versões:

\begin{haskellcode}
-- Versão A
lerA :: IO [Integer]
lerA = do
  x <- readLn
  if x == 0
    then return []
    else do
      xs <- lerA
      return (x:xs)

-- Versão B
lerB :: [Integer] -> IO [Integer]
lerB xs = do
  x <- readLn
  if x == 0
    then return (reverse xs)
    else lerB (x:xs)
\end{haskellcode}

Qual alternativa descreve corretamente a diferença entre elas?

\begin{enumerate}
  \item A versão B não precisa de caso base porque o acumulador substitui a condição de parada.

  \item As duas versões são idênticas quanto à ordem de construção, pois ambas usam o operador \haskellinline|:| no mesmo ponto.

  \item A versão A inclui o sentinela na lista final, enquanto a versão B o descarta.

  \item A versão A reconstrói a lista ao retornar das chamadas recursivas; a versão B acumula na descida e inverte a lista ao final.

  \item A versão A devolve uma lista em ordem inversa e a versão B devolve sempre a lista vazia.


\end{enumerate}

\questionTitle{39}{Comprimentos diferentes na correção de provas}
\begin{flushright}
  \footnotesize
  \textit{\textbf{10: Ações de E/S recursivas}}\\
  \textit{c10-bordas-robustez-05}\\

\end{flushright}
Considere uma função pura auxiliar para corrigir uma prova de múltipla escolha:

\begin{haskellcode}
corrige :: String -> String -> Int
corrige [] [] = 0
corrige (g:gs) (r:rs)
  | g == r    = 1 + corrige gs rs
  | otherwise = corrige gs rs
\end{haskellcode}

Qual problema pode ocorrer se o gabarito e as respostas tiverem comprimentos diferentes?

\begin{enumerate}
  \item A função sempre considera respostas ausentes como corretas, pois \haskellinline|[]| vale \haskellinline|True|.

  \item A função inverte automaticamente a maior lista para tentar alinhar os elementos.

  \item A função passa a ter tipo \haskellinline|IO Int|, pois listas de tamanhos diferentes exigem entrada e saída.

  \item O compilador rejeita a definição antes da execução, pois não é permitido casar padrões com duas listas.

  \item A função pode falhar por padrões não exaustivos, pois não há casos para \haskellinline|[]| contra lista não vazia ou lista não vazia contra \haskellinline|[]|.


\end{enumerate}

\questionTitle{40}{Comparação conceitual entre let e where}
\begin{flushright}
  \footnotesize
  \textit{\textbf{3: Expressões e definições}}\\
  \textit{c03-let-escopo-04}\\

\end{flushright}
Qual alternativa apresenta uma diferença conceitual correta entre \haskellinline|let| e \haskellinline|where|?

\begin{enumerate}
  \item \haskellinline|where| sempre tem escopo global, e \haskellinline|let| sempre tem escopo de módulo.

  \item \haskellinline|let ... in ...| é uma expressão e pode aparecer onde expressões são permitidas; \haskellinline|where| é uma cláusula anexada a uma equação.

  \item \haskellinline|let| permite apenas uma definição local, enquanto \haskellinline|where| permite apenas literais numéricos.

  \item Não há diferença: \haskellinline|let| e \haskellinline|where| são apenas grafias alternativas para a mesma construção sintática.

  \item \haskellinline|where| é uma expressão, enquanto \haskellinline|let| só pode aparecer no fim de uma função.


\end{enumerate}

\questionTitle{41}{Avaliação com cálculo local}
\begin{flushright}
  \footnotesize
  \textit{\textbf{5: Expressão condicional}}\\
  \textit{c05-where-avaliacao-02}\\

\end{flushright}
Considere a função:
\begin{haskellcode}
analisa y
  | x > 10    = "Alto"
  | otherwise = "Baixo"
  where
    x = y * y + 1
\end{haskellcode}
Qual é o resultado de \haskellinline|analisa 3|?

\begin{enumerate}
  \item \haskellinline|3|

  \item \haskellinline|10|

  \item \haskellinline|"Alto"|

  \item \haskellinline|"Baixo"|

  \item Ocorre um erro de tipo em \haskellinline|x > 10|.


\end{enumerate}

\questionTitle{42}{Papel dos parênteses em aplicação de função}
\begin{flushright}
  \footnotesize
  \textit{\textbf{3: Expressões e definições}}\\
  \textit{c03-precedencia-associatividade-04}\\

\end{flushright}
Em Haskell, parênteses agrupam expressões, mas não fazem parte da sintaxe obrigatória de chamada de função. Qual alternativa é equivalente a \haskellinline|f (g x) y|?

\begin{enumerate}
  \item \haskellinline|(f (g x)) y|.

  \item \haskellinline|f (g (x y))|.

  \item \haskellinline|(f g) (x y)|.

  \item \haskellinline|f g x y|.

  \item \haskellinline|f ((g x) y)|.


\end{enumerate}

\questionTitle{43}{Análise de assinatura de tipo}
\begin{flushright}
  \footnotesize
  \textit{\textbf{4: Tipos de Dados}}\\
  \textit{c04-assinatura-analise-01}\\

\end{flushright}
Qual das seguintes definições de função corresponde à assinatura de tipo \haskellinline{String -> Int -> Bool}?

\begin{enumerate}
  \item \begin{haskellcode}
verifica b s = not b && null s
\end{haskellcode}

  \item \begin{haskellcode}
verifica s n = length s > n
\end{haskellcode}

  \item \begin{haskellcode}
verifica s n = s ++ show n
\end{haskellcode}

  \item \begin{haskellcode}
verifica n s = n > 0 && s == "OK"
\end{haskellcode}

  \item \begin{haskellcode}
verifica s = s == "OK"
\end{haskellcode}


\end{enumerate}

\questionTitle{44}{Literais fracionários e fromRational}
\begin{flushright}
  \footnotesize
  \textit{\textbf{7: Polimorfismo}}\\
  \textit{cap07-literais-fromrational-contexto}\\

\end{flushright}
Qual afirmação descreve corretamente literais fracionários como \haskellinline{3.14}?
\begin{enumerate}
  \item Eles exigem \haskellinline{Integral a}, pois possuem uma representação decimal.
  \item Eles têm sempre tipo \haskellinline{Double}, independentemente do contexto.
  \item Eles são interpretados por meio de \haskellinline{fromRational} e recebem inicialmente uma restrição \haskellinline{Fractional a => a}.
  \item Eles podem ser usados diretamente como \haskellinline{Int}, pois Haskell arredonda automaticamente.
  \item Eles são convertidos por \haskellinline{fromInteger}, assim como \haskellinline{42}.

\end{enumerate}

\questionTitle{45}{A ideia de mundo em IO}
\begin{flushright}
  \footnotesize
  \textit{\textbf{8: Programas interativos}}\\
  \textit{c08-modelo-io-05}\\

\end{flushright}
No modelo conceitual apresentado para programas interativos em Haskell, qual é a melhor interpretação da ideia de \emph{mundo}?

\begin{enumerate}
  \item É uma estrutura de dados concreta do Prelude que o programador deve passar manualmente para cada função.

  \item É apenas outro nome para o tipo \haskellinline|String| usado nas entradas do usuário.

  \item É um módulo obrigatório que precisa ser importado para usar \haskellinline|getLine| e \haskellinline|putStrLn|.

  \item É uma forma conceitual de representar o ambiente externo com o qual o programa interage, como terminal, arquivos e outros dispositivos, preservando a separação entre valores puros e ações de E/S.

  \item É uma variável global chamada \haskellinline|World| que pode ser lida e modificada por qualquer expressão pura.


\end{enumerate}

\questionTitle{46}{Definição de linguagens multiparadigma}
\begin{flushright}
  \footnotesize
  \textit{\textbf{1: Paradigmas de programação}}\\
  \textit{c01-conceitos-paradigmas-04}\\

\end{flushright}
Linguagens modernas de amplo uso na indústria, como Python e C++, são frequentemente classificadas como \textbf{multiparadigma}. O que esta classificação indica do ponto de vista da arquitetura de software?

\begin{enumerate}
  \item Indica que a linguagem é capaz de compilar o mesmo código-fonte para múltiplas arquiteturas de hardware diferentes (como x86, ARM e RISC-V).

  \item Indica que o compilador da linguagem foi escrito utilizando um paradigma distinto daquele que é exposto aos programadores que a utilizam.

  \item Indica que a linguagem possui simultaneamente um interpretador em tempo real e um compilador \emph{Ahead-of-Time} (AOT).

  \item Indica que a linguagem exige que o desenvolvedor escolha um único paradigma no momento da criação do projeto, desabilitando permanentemente os comandos dos demais paradigmas.

  \item Indica que a linguagem fornece suporte nativo a técnicas e abstrações de múltiplos paradigmas (como o estruturado, o orientado a objetos e o funcional), permitindo ao desenvolvedor mesclar essas abordagens.


\end{enumerate}

\questionTitle{47}{Identificação de forma normal}
\begin{flushright}
  \footnotesize
  \textit{\textbf{3: Expressões e definições}}\\
  \textit{c03-layout-avaliacao-04}\\

\end{flushright}
Qual das alternativas abaixo está em forma normal no sentido apresentado no capítulo, isto é, não precisa de mais reduções para ser um valor?

\begin{enumerate}
  \item \haskellinline|(\x -> x) 5|.

  \item \haskellinline|let x = 3 in x + 1|.

  \item \haskellinline|7|.

  \item \haskellinline|sqrt 16|.

  \item \haskellinline|1 + 2 * 3|.


\end{enumerate}

\questionTitle{48}{Fibonacci com acumuladores}
\begin{flushright}
  \footnotesize
  \textit{\textbf{9: Funções recursivas}}\\
  \textit{c09-cauda-fibonacci-01}\\

\end{flushright}
Em uma versão de Fibonacci com acumuladores, mantém-se um par \haskellinline|(a,b)| de números consecutivos da sequência. Qual atualização preserva corretamente essa interpretação?

\begin{enumerate}
  \item Atualizar \haskellinline|(a,b)| para \haskellinline|(a-1,b-1)|.

  \item Atualizar \haskellinline|(a,b)| para \haskellinline|(a,b+1)|.

  \item Atualizar \haskellinline|(a,b)| para \haskellinline|(b,a+b)|.

  \item Atualizar \haskellinline|(a,b)| para \haskellinline|(b,a*b)|.

  \item Atualizar \haskellinline|(a,b)| para \haskellinline|(a+b,b)|.


\end{enumerate}

\questionTitle{49}{Equações como definições}
\begin{flushright}
  \footnotesize
  \textit{\textbf{3: Expressões e definições}}\\
  \textit{c03-ghci-definicoes-04}\\

\end{flushright}
Analise a definição:

\begin{haskellcode}
quadrado x = x * x
\end{haskellcode}

Qual alternativa descreve corretamente sua estrutura?

\begin{enumerate}
  \item A definição é inválida porque toda função Haskell deve ter seus parâmetros entre parênteses.

  \item \haskellinline|quadrado| é o nome definido, \haskellinline|x| é um parâmetro formal e \haskellinline|x * x| é a expressão do lado direito da equação.

  \item \haskellinline|quadrado| é um construtor, pois aparece à esquerda do sinal de igualdade.

  \item \haskellinline|quadrado x| é uma tupla formada por dois nomes, e \haskellinline|x * x| é um comentário.

  \item \haskellinline|x * x| é avaliado imediatamente para um único número no momento em que o arquivo é carregado.


\end{enumerate}

\questionTitle{50}{Homogeneidade de listas}
\begin{flushright}
  \footnotesize
  \textit{\textbf{4: Tipos de Dados}}\\
  \textit{c04-checagem-listas-01}\\

\end{flushright}
Em Haskell, se um programador tentar avaliar a expressão \haskellinline{[1, True, "texto"]}, o que acontecerá?

\begin{enumerate}
  \item Ocorrerá um erro em tempo de compilação, pois todos os elementos de uma lista devem ter o mesmo tipo.

  \item O compilador criará uma tupla automaticamente.

  \item Ocorrerá um erro apenas em tempo de execução, ao tentar usar os elementos.

  \item Os valores serão implicitamente convertidos para \haskellinline{String}.

  \item A expressão será aceita e o tipo inferido será \haskellinline{Any}.


\end{enumerate}

\questionTitle{51}{Associatividade da seta de tipo}
\begin{flushright}
  \footnotesize
  \textit{\textbf{4: Tipos de Dados}}\\
  \textit{c04-funcao-associatividade-01}\\

\end{flushright}
Devido à associatividade à direita da seta de tipo (\haskellinline{->}), como a assinatura \haskellinline{Int -> Bool -> String} é implicitamente parentizada?

\begin{enumerate}
  \item \haskellinline{(Int -> Bool) -> String}

  \item A seta não possui associatividade definida.

  \item \haskellinline{Int -> (Bool -> String)}

  \item \haskellinline{Int -> Bool -> (String)}

  \item \haskellinline{(Int, Bool) -> String}


\end{enumerate}

\questionTitle{52}{Identificação de violação de pureza (I/O temporal)}
\begin{flushright}
  \footnotesize
  \textit{\textbf{1: Paradigmas de programação}}\\
  \textit{c01-transparencia-referencial-02}\\

\end{flushright}
Considere uma função \haskellinline|obterHoraAtual| que consulta o relógio interno do sistema operacional e retorna uma representação do momento exato em que foi chamada.

Por que a existência desta função viola a propriedade de pureza em um modelo matemático estrito?

\begin{enumerate}
  \item Ela não viola a pureza; contanto que a função não modifique a hora do sistema (apenas efetue uma leitura), ela continua sendo referencialmente transparente.

  \item Porque ela consome excessivos ciclos de processamento do \emph{garbage collector}, atrasando a execução do fluxo principal da aplicação.

  \item Porque seu valor de retorno depende de um estado global em mutação constante fora do controle do programa, fazendo com que chamadas com os mesmos argumentos produzam resultados diferentes.

  \item Porque a leitura do relógio do sistema exige chamadas de baixo nível em \emph{Assembly}, que quebram as abstrações do compilador Haskell.

  \item Porque o tempo é uma grandeza contínua e linguagens funcionais só conseguem modelar estruturas de dados e matrizes estritamente discretas.


\end{enumerate}

\questionTitle{53}{Identificadores simbólicos}
\begin{flushright}
  \footnotesize
  \textit{\textbf{3: Expressões e definições}}\\
  \textit{c03-arquivos-identificadores-04}\\

\end{flushright}
Qual alternativa apresenta um identificador simbólico apropriado para definir um novo operador infixo comum, sem usar a forma reservada a construtores?

\begin{enumerate}
  \item \haskellinline|op+|.

  \item \haskellinline|A++|.

  \item \haskellinline|:++|.

  \item \haskellinline|let|.

  \item \haskellinline|<++>|.


\end{enumerate}

\questionTitle{54}{Equivalência com True}
\begin{flushright}
  \footnotesize
  \textit{\textbf{5: Expressão condicional}}\\
  \textit{c05-otherwise-02}\\

\end{flushright}
Qual definição é equivalente ao uso usual de \haskellinline|otherwise| na última guarda?

\begin{enumerate}
  \item Substituir \haskellinline|otherwise| por \haskellinline|False|.

  \item Substituir \haskellinline|otherwise| por \haskellinline|True|.

  \item Substituir \haskellinline|otherwise| por \haskellinline|else|.

  \item Substituir \haskellinline|otherwise| por \haskellinline|()|.

  \item Substituir \haskellinline|otherwise| por \haskellinline|undefined|.


\end{enumerate}

\questionTitle{55}{Uso infixo de uma função prefixa}
\begin{flushright}
  \footnotesize
  \textit{\textbf{3: Expressões e definições}}\\
  \textit{c03-aplicacao-notacoes-04}\\

\end{flushright}
Qual alternativa é equivalente à expressão \haskellinline|div (sum [1,2,3,4]) 3| usando a função \haskellinline|div| em notação infixa?

\begin{enumerate}
  \item \haskellinline|div `sum [1,2,3,4]` 3|.

  \item \haskellinline|sum [1,2,3,4] div 3|.

  \item \haskellinline|(sum [1,2,3,4], 3) `div`|.

  \item \haskellinline|3 `div` (sum [1,2,3,4])|.

  \item \haskellinline|(sum [1,2,3,4]) `div` 3|.


\end{enumerate}

\questionTitle{56}{getLine seguido de read}
\begin{flushright}
  \footnotesize
  \textit{\textbf{8: Programas interativos}}\\
  \textit{c08-conversao-validacao-05}\\

\end{flushright}
Qual trecho expressa corretamente a ideia de ler uma linha e convertê-la para \haskellinline|Int|?

\begin{enumerate}
  \item \begin{minted}{haskell}
let n <- getLine :: Int
\end{minted}

  \item \begin{minted}{haskell}
linha <- getLine
let n = read linha :: Int
\end{minted}

  \item \begin{minted}{haskell}
linha <- read getLine
let n = linha :: Int
\end{minted}

  \item \begin{minted}{haskell}
n <- read getLine :: IO Int
\end{minted}

  \item \begin{minted}{haskell}
let linha = getLine
let n = read linha :: Int
\end{minted}


\end{enumerate}

\questionTitle{57}{Problema menor do mesmo tipo}
\begin{flushright}
  \footnotesize
  \textit{\textbf{9: Funções recursivas}}\\
  \textit{c09-fundamentos-recursao-estrutural-01}\\

\end{flushright}
Em uma definição recursiva de \haskellinline|potDois|, a chamada recursiva principal é
\haskellinline|potDois (n - 1)| e o caso base é \haskellinline|n == 0|. Qual é a melhor interpretação dessa chamada?

\begin{enumerate}
  \item Ela preserva o tipo de problema, mas o torna menor porque subtrai uma unidade de n, aproximando a computação do caso base.

  \item Ela é equivalente a uma estrutura \haskellinline|while| imperativa escrita diretamente em Haskell.

  \item Ela impede a terminação, pois toda chamada recursiva necessariamente se afasta do caso base.

  \item Ela muda o problema para outro tipo de cálculo, por isso deixa de ser uma definição recursiva.

  \item Ela é chamada apenas depois que o caso base já foi alcançado.


\end{enumerate}

\questionTitle{58}{Quantidade de termos na aproximação de pi}
\begin{flushright}
  \footnotesize
  \textit{\textbf{10: Ações de E/S recursivas}}\\
  \textit{c10-rastreamento-05}\\

\end{flushright}
Considere a versão revisada da aproximação de \(\pi\):

\begin{haskellcode}
valorPi :: Integer -> Double
valorPi n
  | n <= 0    = 0
  | otherwise = 4 * soma 0 (n - 1)

soma :: Integer -> Integer -> Double
soma i n
  | i <= n    = (-1) ^ i / (2 * fromIntegral i + 1)
                  + soma (i + 1) n
  | otherwise = 0
\end{haskellcode}

Se \haskellinline|valorPi 4| é avaliada, quais valores de \haskellinline|i| são usados nos termos da série?

\begin{enumerate}
  \item \haskellinline|1, 2, 3, 4|, pois o parâmetro indica o maior índice.

  \item \haskellinline|4| apenas, pois \haskellinline|soma| avalia somente o último termo.

  \item \haskellinline|0, 1, 2, 3|, pois o parâmetro indica a quantidade de termos e o limite superior é \haskellinline|n - 1|.

  \item \haskellinline|0, 1, 2, 3, 4|, pois a função sempre usa \haskellinline|n + 1| termos.

  \item Nenhum valor, pois \haskellinline|valorPi 4| cai no caso \haskellinline|n <= 0|.


\end{enumerate}

\questionTitle{59}{Reuso uniforme e comportamento por instância}
\begin{flushright}
  \footnotesize
  \textit{\textbf{7: Polimorfismo}}\\
  \textit{cap07-comparacao-reuso-e-instancias}\\

\end{flushright}
Considere as funções:

\begin{minted}{haskell}
length :: [a] -> Int
show   :: Show a => a -> String
\end{minted}

Qual comparação é correta?

\begin{enumerate}
  \item As duas são paramétricas puras, pois ambas usam a variável de tipo \haskellinline{a}.
  \item \haskellinline{show} não usa polimorfismo, pois sempre retorna \haskellinline{String}.
  \item \haskellinline{length} é uniformemente aplicável a listas de qualquer tipo de elemento; \haskellinline{show} depende de uma instância \haskellinline{Show} para o tipo do valor.
  \item \haskellinline{length} depende de \haskellinline{Show a}, pois precisa imprimir cada elemento antes de contar.
  \item As duas são monomórficas, pois retornam tipos concretos.

\end{enumerate}

\questionTitle{60}{A Crise do Software e a solução funcional}
\begin{flushright}
  \footnotesize
  \textit{\textbf{1: Paradigmas de programação}}\\
  \textit{c01-engenharia-software-02}\\

\end{flushright}
O termo \emph{Crise do Software} surgiu historicamente devido à dificuldade de entregar sistemas complexos no prazo, dentro do orçamento e isentos de falhas catastróficas em tempo de execução. Como o paradigma funcional atua para mitigar diretamente a raiz estrutural dessa crise?

\begin{enumerate}
  \item Reduzindo o tempo de desenvolvimento ao automatizar a geração de código através de modelos probabilísticos baseados em inteligência artificial no sistema de tipos dinâmicos.

  \item Elevando o nível de abstração focado na declaração (o que fazer) e adotando bases que eliminam o estado mutável não documentado, o que reduz drasticamente a complexidade do rastreamento de defeitos.

  \item Eliminando de forma definitiva a necessidade de escrever casos de teste unitário, já que o processo de compilação cruzada garante a integridade incondicional das regras de negócio.

  \item Restringindo rigorosamente a arquitetura dos projetos a um limite máximo de linhas de código por pacote para evitar que os softwares atinjam uma escala insustentável.

  \item Forçando a adoção de metodologias de desenvolvimento ágil (\emph{Agile/Scrum}) através de ferramentas nativas embutidas diretamente nos compiladores das linguagens funcionais.


\end{enumerate}

\questionTitle{61}{Rastreamento de Fibonacci}
\begin{flushright}
  \footnotesize
  \textit{\textbf{9: Funções recursivas}}\\
  \textit{c09-formas-fibonacci-rastreamento-01}\\

\end{flushright}
Usando a definição do capítulo,
\begin{haskellcode}
fib n
  | n == 0 = 0
  | n == 1 = 1
  | n > 1  = fib (n - 2) + fib (n - 1)
\end{haskellcode}
qual é o valor de \haskellinline|fib 5|?

\begin{enumerate}
  \item \haskellinline|8|

  \item \haskellinline|4|

  \item \haskellinline|5|

  \item \haskellinline|3|

  \item \haskellinline|10|


\end{enumerate}

\questionTitle{62}{Remoção de módulos do escopo}
\begin{flushright}
  \footnotesize
  \textit{\textbf{2: Ambiente de desenvolvimento}}\\
  \textit{c02-ghci-modulos-04}\\

\end{flushright}
O atalho \texttt{:m} (abreviação de \texttt{:module}) do GHCi oferece uma flexibilidade adicional na gestão do escopo que a palavra-chave tradicional \haskellinline{import} não possui no terminal. Qual é essa flexibilidade?

\begin{enumerate}
  \item O comando \texttt{:m} força a avaliação estrita de toda a biblioteca sem utilizar as otimizações preguiçosas.

  \item O comando \texttt{:m} executa as funções do módulo importado em \emph{threads} separadas visando alto paralelismo.

  \item O comando \texttt{:m} realiza um \emph{dump} da memória do GHCi e grava o ambiente ativo em um novo arquivo \texttt{.hs}.

  \item O comando \texttt{:m} contorna a necessidade de conectividade com a internet ao buscar módulos indisponíveis no cache.

  \item O comando \texttt{:m} permite a utilização do prefixo de subtração (\texttt{-}), como em \texttt{:m -Data.List}, propiciando a remoção rápida de módulos do escopo interativo atual.


\end{enumerate}

\questionTitle{63}{Diagnóstico de ordem aparente}
\begin{flushright}
  \footnotesize
  \textit{\textbf{8: Programas interativos}}\\
  \textit{c08-buffering-05}\\

\end{flushright}
Um aluno observa que o terminal parece aguardar a entrada antes de mostrar o \emph{prompt}. Qual diagnóstico é mais apropriado para o seguinte padrão?

\begin{minted}{haskell}
putStr "Digite o valor: "
valor <- getLine
\end{minted}

\begin{enumerate}
  \item Haskell executou \haskellinline|getLine| antes de \haskellinline|putStr| por causa da avaliação preguiçosa.

  \item A ação \haskellinline|putStr| não pode ser usada antes de \haskellinline|getLine|.

  \item O programa está correto apenas se \haskellinline|valor| tiver tipo \haskellinline|IO String|.

  \item A variável \haskellinline|valor| deveria ser definida com \haskellinline|let|.

  \item A ordem das ações está correta; o problema provável é a saída bufferizada, que pode atrasar a exibição do texto produzido por \haskellinline|putStr|.


\end{enumerate}

\questionTitle{64}{Inferência de restrições implícitas}
\begin{flushright}
  \footnotesize
  \textit{\textbf{7: Polimorfismo}}\\
  \textit{cap07-inferencia-classes-implicitas}\\

\end{flushright}
Considere a função:

\begin{minted}{haskell}
positivoMaiorQue x y = x > y && y > 0
\end{minted}

Qual é a assinatura mais geral?

\begin{enumerate}
  \item \mintinline{haskell}{positivoMaiorQue :: Num a => a -> a -> Bool}
  \item \mintinline{haskell}{positivoMaiorQue :: Ord a => a -> a -> Bool}
  \item \mintinline{haskell}{positivoMaiorQue :: Integral a => a -> a -> Bool}
  \item \mintinline{haskell}{positivoMaiorQue :: Bool -> Bool -> Bool}
  \item \mintinline{haskell}{positivoMaiorQue :: (Ord a, Num a) => a -> a -> Bool}

\end{enumerate}

\questionTitle{65}{Detalhes profundos com o comando :info}
\begin{flushright}
  \footnotesize
  \textit{\textbf{2: Ambiente de desenvolvimento}}\\
  \textit{c02-ghci-inspecao-04}\\

\end{flushright}
Considere que o aluno já sabe o tipo do operador de adição (\haskellinline{+}), mas deseja investigar em qual módulo interno ele foi definido e qual é o nível matemático de sua precedência (\emph{fixity}). Qual comando do GHCi fornece essa camada extra de metadados?

\begin{enumerate}
  \item O comando \texttt{:help} (ou \texttt{:?}).

  \item O comando \texttt{:type} (ou \texttt{:t}).

  \item O comando \texttt{:info} (ou \texttt{:i}).

  \item O comando \texttt{:browse} (ou \texttt{:b}).

  \item O comando \texttt{:show} (ou \texttt{:s}).


\end{enumerate}

\questionTitle{66}{Cabeça, cauda e construção}
\begin{flushright}
  \footnotesize
  \textit{\textbf{6: Estruturas de dados básicas}}\\
  \textit{c06-listas-estrutura-02}\\

\end{flushright}
Qual é o resultado da expressão abaixo?
\begin{minted}{haskell}
tail ("a" ++ "bc")
\end{minted}

\begin{enumerate}
  \item \haskellinline{ "c" }

  \item \haskellinline{ "a" }

  \item \haskellinline{ "abc" }

  \item A expressão resulta em erro de execução.

  \item \haskellinline{ "bc" }


\end{enumerate}

\questionTitle{67}{String como lista de Char}
\begin{flushright}
  \footnotesize
  \textit{\textbf{4: Tipos de Dados}}\\
  \textit{c04-tipos-string-01}\\

\end{flushright}
Sabendo que \haskellinline{String} é um tipo sinônimo para \haskellinline{[Char]}, qual das seguintes expressões é \textbf{equivalente} a \haskellinline{"FP"}?

\begin{enumerate}
  \item \haskellinline{["F", "P"]}

  \item \haskellinline{'FP'}

  \item \haskellinline{('F', 'P')}

  \item \haskellinline{"F" : "P"}

  \item \haskellinline{['F', 'P']}


\end{enumerate}

\questionTitle{68}{getChar e getLine}
\begin{flushright}
  \footnotesize
  \textit{\textbf{8: Programas interativos}}\\
  \textit{c08-entrada-saida-05}\\

\end{flushright}
Qual alternativa compara corretamente \haskellinline|getChar| e \haskellinline|getLine|?

\begin{enumerate}
  \item Ambas têm tipo \haskellinline|IO String|, mas \haskellinline|getChar| retorna uma string de tamanho 1.

  \item \haskellinline|getChar| é pura; \haskellinline|getLine| é uma ação de E/S.

  \item \haskellinline|getChar| tem tipo \haskellinline|IO Char| e lê um caractere; \haskellinline|getLine| tem tipo \haskellinline|IO String| e lê uma linha de texto.

  \item \haskellinline|getChar| imprime um caractere; \haskellinline|getLine| imprime uma linha.

  \item \haskellinline|getLine| recebe uma string de prompt como argumento, enquanto \haskellinline|getChar| não recebe argumentos.


\end{enumerate}

\questionTitle{69}{Diagnóstico: Erro de diretório no load}
\begin{flushright}
  \footnotesize
  \textit{\textbf{2: Ambiente de desenvolvimento}}\\
  \textit{c02-fluxo-trabalho-04}\\

\end{flushright}
Um aluno salva o seu código sob o nome \texttt{Pratica.hs} na pasta \emph{Documentos} e, na sequência, abre um terminal recém-iniciado e digita o comando \texttt{:load Pratica.hs} no GHCi. O interpretador retorna o erro de sistema \emph{Target "Pratica.hs" is not a module name or a source file}. 

Qual é o passo crucial que foi omitido no fluxo de inicialização do ambiente?

\begin{enumerate}
  \item O GHCi obriga que o arquivo possua a extensão binária \texttt{.bin} ao invés da extensão de texto plano \texttt{.hs}.

  \item O aluno não possuía a licença corporativa do \emph{Haskell Language Server} ativada para permitir a leitura remota de discos.

  \item O aluno deveria ter forçado a importação global com o comando de sistema \texttt{:sudo load Pratica.hs}.

  \item O aluno esqueceu de garantir que o terminal do sistema operacional estava no mesmo diretório (pasta) onde o arquivo foi salvo antes de chamar o executável \texttt{ghci}.

  \item O aluno nomeou o arquivo com a primeira letra em maiúscula (\texttt{P}), o que é estritamente proibido pela especificação sintática de módulos.


\end{enumerate}


\end{multicols}
\makeanswersheet{UFOP}{BCC222: Programação Funcional}{José Romildo}{2026/1}{Primeira Prova}{07}{69}{25}{7}
\gabarito{07}{D,E,C,C,C,B,D,E,A,A,D,E,C,D,B,E,B,C,D,D,B,B,A,E,A,D,E,B,E,E,E,B,D,D,B,A,D,D,E,B,D,A,B,C,D,E,C,C,B,A,C,C,E,B,E,B,A,C,C,B,C,E,E,E,C,E,E,C,D}
\end{document}
